% ============================================================
%  Paper A — Supplementary Information
% ============================================================

% Reset figure and table counters + prefixes so SI elements are
% auto-numbered "Supplementary Figure S1, S2, ..." and
% "Supplementary Table S1, S2, ...". Without these the ED figure counter
% would bleed through (captions would read "Extended Data Fig. 10, 11,
% ...") and the global table counter would give "Table 1, 2, 3" on our
% SI tables.
\setcounter{figure}{0}
\setcounter{table}{0}
\renewcommand{\figurename}{Supplementary Figure}
\renewcommand{\tablename}{Supplementary Table}
\renewcommand{\thefigure}{S\arabic{figure}}
\renewcommand{\thetable}{S\arabic{table}}

% Auto-numbered Supplementary Notes (narrative analyses), numbered 1..N
% in order of appearance and cross-referenced via \ref. Figure/table-holder
% subsections keep descriptive headers; they are identified by their
% Supplementary Figure/Table number.
\newcounter{snote}
\newcommand{\suppnote}[2]{%
  \refstepcounter{snote}%
  \subsection*{Supplementary Note~\arabic{snote}: #1}%
  \label{#2}%
}

% ── GPU-derived supplementary figures ──────────────────────────────────

\subsection*{Organism growth rates (classified corpus)}

\begin{figure}[H]
  \centering
  \includegraphics[width=0.72\textwidth]{S1_growth_rates}
  \caption{\textbf{Annual growth rates by organism (zero-shot classified corpus,
    2.77 million papers).}
    CAGR computed over the 2010s and 2020s for each organism category.
    Rice leads crop-organism growth in the 2010s (CAGR 4.9\%); most organism
    groups show slower or negative growth in the 2020s consistent with
    field-level plateau.}
  \label{fig:S1}
\end{figure}

\subsection*{Paradigm classification and citation impact}

\begin{figure}[H]
  \centering
  \includegraphics[width=0.72\textwidth]{S2_paradigm_citations}
  \caption{\textbf{Citation impact by paradigm and organism.}
    Box plots of five-year forward citation counts for applied versus fundamental
    papers, stratified by organism.
    Applied \textit{Arabidopsis} papers show the highest median citation counts;
    fundamental papers in crop organisms show high variance reflecting a small
    number of highly cited mechanistic studies.}
  \label{fig:S2}
\end{figure}

\subsection*{Topic diversity by year}

\begin{figure}[H]
  \centering
  \includegraphics[width=0.72\textwidth]{S3_topic_diversity}
  \caption{\textbf{Shannon entropy of macro-theme distribution by year
    (1990--2024).}
    Entropy declined from 1.82 bits (1990s mean) to 1.14 bits (2020s mean),
    reflecting concentration of the research agenda around a smaller set of
    dominant themes.
    Shaded region: 95\% bootstrap confidence interval.}
  \label{fig:S3}
\end{figure}

\subsection*{Disruption index by organism}

\begin{figure}[H]
  \centering
  \includegraphics[width=0.72\textwidth]{S6_disruption_index}
  \caption{\textbf{CD disruption index by organism, decade resolution.}
    \textit{Arabidopsis} has the most negative mean CD index ($-0.001$),
    indicating consolidation; crop species (maize 0.095, soybean 0.102,
    cotton 0.099) are consistently more disruptive.
    Error bars: standard error of the mean across decades.
    Kruskal--Wallis test: $H = 12.9$, $p = 0.17$, $n_{\mathrm{groups}} = 10$.}
  \label{fig:S6}
\end{figure}

\subsection*{Cross-citation matrix between organism categories}

\begin{figure}[H]
  \centering
  \includegraphics[width=0.72\textwidth]{S9_cross_citation}
  \caption{\textbf{Cross-citation matrix between organism categories
    (2020--2024).}
    Cell values give the fraction of citations from organism~$i$ (row) to
    papers in organism category~$j$ (column).
    \textit{Arabidopsis} has the highest self-citation rate (61.6\%);
    crop organisms cite \textit{Arabidopsis} at rates of 10--20\%,
    confirming continued intellectual dependence on the model organism.}
  \label{fig:S9}
\end{figure}

\subsection*{International collaboration by organism}

\begin{figure}[H]
  \centering
  \includegraphics[width=0.72\textwidth]{S10_collaboration}
  \caption{\textbf{International collaboration rate by organism (2020--2024).}
    \textit{Arabidopsis} has the highest international collaboration rate
    (18.9\%), followed by other model organisms (16.2\%) and barley (16.1\%).
    Soybean (11.3\%) and non-specific papers (9.8\%) have the lowest rates.}
  \label{fig:S10}
\end{figure}

% ── Traditional supplementary tables ───────────────────────────────────

\subsection*{Orphan crops ranked by research-attention gap}

Plants are ranked by attention gap index (log$_2$ ratio of papers per trillion
kilocalories). Negative values indicate under-studied species relative to their
caloric importance. Data from FAO production statistics and OpenAlex corpus
(1900--2024).

\begin{table}[H]
\centering
\caption{Crops ranked by research-attention gap (most under-studied first).}
\label{tab:S1}
\begin{tabular}{rlrrrr}
\toprule
Rank & Crop & Publications & Production (Mt) & Papers/Mt & Attention gap \\
\midrule
1  & Sugarcane       & 30,220  & 1,900 & 15.9    & $-4.27$ \\
2  & Cassava         & 21,135  & 310   & 68.2    & $-2.17$ \\
3  & Maize           & 138,597 & 1,160 & 119.5   & $-1.36$ \\
4  & Potato          & 52,825  & 375   & 140.9   & $-1.12$ \\
5  & Wheat           & 152,000 & 780   & 194.9   & $-0.65$ \\
6  & Plantain        & 29,912  & 120   & 249.3   & $-0.30$ \\
7  & Soybean         & 90,199  & 350   & 257.7   & $-0.25$ \\
8  & Rice            & 168,437 & 520   & 323.9   & $+0.08$ \\
9  & Barley          & 52,570  & 155   & 339.2   & $+0.15$ \\
10 & Tomato          & 79,974  & 187   & 427.7   & $+0.48$ \\
11 & Sorghum         & 33,419  & 60    & 557.0   & $+0.86$ \\
12 & Millet          & 20,263  & 30    & 675.4   & $+1.14$ \\
13 & Tobacco         & 46,616  & 6.3   & 7,399.4 & $+4.59$ \\
\bottomrule
\end{tabular}
\end{table}

\bigskip

\subsection*{Top sleeping-beauty papers on minor crops}

Papers with highest beauty index ($B$) from the sleeping-beauty analysis
using the Ke et al.\ (2015) metric. Dormancy defined as years from
publication to awakening (first year of sustained citation acceleration).

\begin{table}[H]
\centering
\caption{Top 10 sleeping beauty papers by $B$-index.}
\label{tab:S2}
\begin{tabular}{rlrlrr}
\toprule
Rank & Crop & Year & Title (abbreviated) & Citations & $B$-index \\
\midrule
1  & Moringa   & 1980 & Drumstick: A multipurpose Indian vegetable            & 580   & 83.1 \\
2  & Moringa   & 2006 & \textit{M.\ oleifera}: multiple medicinal uses        & 1,600 & 63.0 \\
3  & Moringa   & 2005 & Review of medical evidence, Part 1                    & ---   & 55.2 \\
4  & Moringa   & 1996 & Nutritional value of whole \& extracted leaves         & 539   & 45.8 \\
5  & Moringa   & 2003 & Antioxidant properties of drumstick tree               & 1,382 & 43.7 \\
6  & Moringa   & 1991 & Horseradish tree: a boon to arid lands                & 449   & 42.5 \\
7  & Moringa   & 2001 & Potential for agricultural \& industrial uses          & 458   & 41.9 \\
8  & Millet    & 1982 & Peroxidase in excised ragi leaves                     & 182   & 39.8 \\
9  & Moringa   & 1999 & The miracle tree: natural nutrition for tropics        & 416   & 38.5 \\
10 & Amaranth  & 1987 & Planting date influence on Palmer amaranth growth      & 235   & 28.8 \\
\bottomrule
\end{tabular}
\end{table}

\bigskip

\subsection*{Method diffusion parameters}

Fitted logistic S-curve parameters for 13 methodological categories detected
in plant science publications. Half-life is the time from 10\% to 50\% of
carrying capacity $K$.

\begin{table}[H]
\centering
\caption{Logistic diffusion model parameters for methodologies in plant science.}
\label{tab:S3}
\begin{tabular}{lrrrrrrr}
\toprule
Method & Origin & Total papers & $K$ & $r$ (yr$^{-1}$) & $t_{\mathrm{mid}}$ & Half-life (yr) & $R^2$ \\
\midrule
Long-read seq.    & 2014 & 2,272  & 3,127   & 0.512 & 2022.1 & 8.6  & 1.000 \\
GBS/RAD-seq       & 2011 & 1,629  & 1,740   & 0.494 & 2019.1 & 8.9  & 0.999 \\
CRISPR            & 2012 & 4,370  & 5,964   & 0.475 & 2022.0 & 9.2  & 0.999 \\
Pangenomics       & 2015 & 753    & 1,533   & 0.402 & 2024.0 & 10.9 & 0.998 \\
RNA-seq           & 2008 & 13,762 & 16,899  & 0.394 & 2020.3 & 11.1 & 0.999 \\
Machine learning  & 2010 & 27,329 & 163,974 & 0.354 & 2028.6 & 12.4 & 0.995 \\
Microarray        & 1995 & 6,365  & 6,374   & 0.311 & 2011.0 & 14.1 & 0.999 \\
Proteomics        & 1997 & 13,775 & 16,449  & 0.237 & 2017.4 & 18.6 & 0.999 \\
Metabolomics      & 2000 & 9,597  & 31,809  & 0.218 & 2027.9 & 20.2 & 0.999 \\
PCR               & 1985 & 44,077 & 50,283  & 0.177 & 2013.8 & 24.8 & 0.999 \\
Confocal microsc. & 1985 & 3,536  & 4,433   & 0.167 & 2015.9 & 26.2 & 0.999 \\
Flow cytometry    & 1975 & 4,844  & 6,139   & 0.146 & 2015.2 & 30.1 & 1.000 \\
Single-cell       & 2015 & 1,607  & 9,642   & 0.112 & 2039.1 & 39.4 & 0.995 \\
\bottomrule
\end{tabular}
\end{table}

\suppnote{Temporal placebo test (reorganization signal)}{sec:temporal-placebo}

To rule out the possibility that our ``reorganization around climate-relevant
crops'' finding reflects generic publication growth rather than a genuine
structural shift, we applied the same signal detector --- the linear slope of
annual climate-relevant-crop abstract share --- to a placebo window
(1970--1999), in which the climate-discourse framing was not yet established,
and to the treatment window (2000--2024).

The placebo-window slope is $-2.3 \times 10^{-4}$ per year
(95\% bootstrap CI $[-3.7, -0.9] \times 10^{-4}$, 1{,}000 resamples of
year--share pairs), indicating a flat or very slightly declining signal in the
pre-2000 era.
The treatment-window slope is $+4.5 \times 10^{-4}$ per year
(95\% bootstrap CI $[+2.7, +6.2] \times 10^{-4}$), i.e.\ the two windows'
confidence intervals do not overlap.
A two-sided permutation test (1{,}000 shuffles of pooled year--share pairs,
statistic = slope gap) rejects the null of a single underlying slope at
$p = 0.001$.

The signal is era-specific and cannot be reduced to background publication
growth (Figure~\ref{fig:s-temporal-placebo}).

\begin{figure}[H]
  \centering
  \includegraphics[width=0.6\linewidth]{fig_temporal_placebo.pdf}
  \caption{\textbf{Temporal placebo test.} Reorganization effect size (linear
  slope of annual climate-relevant-crop abstract share) for the 1970--1999
  placebo window (grey) and the 2000--2024 treatment window (blue). Error
  bars show 95\% bootstrap confidence intervals from 1{,}000 resamples of
  year--share pairs. A permutation test under the exchangeable-window null
  gives $p = 0.001$ (1{,}000 shuffles).}
  \label{fig:s-temporal-placebo}
\end{figure}

\suppnote{Growth-model selection --- logistic vs Gompertz vs Richards}{sec:growth-model-aic}

We chose the logistic model for the field-growth plateau claim in the
main text. To address the reviewer question ``why logistic over
alternative sigmoidal growth models?'' we fit three candidates to the
deduplicated (works\_clean) annual publication counts 1990--2024
($n = 35$) by non-linear least squares: the three-parameter logistic,
the three-parameter Gompertz, and the four-parameter Richards
(generalised logistic). Logistic and Richards are fit with year as the
time axis (matching the main-text parameterisation); Gompertz is fit on
a translated $t = \mathrm{year}-1990$ axis for numerical stability, with
the carrying-capacity $K$ invariant under the shift.
Table~\ref{tab:S7} reports the parameter fits, the residual sum of
squares (RSS), and the Akaike information criterion corrected for small
samples (AICc); Figure~\ref{fig:s-growth-aic} overlays the three fits on
the observed series.

\begin{table}[H]
\centering
\caption{Three-way comparison of sigmoidal growth-model fits to
plant-science annual publication counts (1990--2024, works\_clean).}
\label{tab:S7}
\begin{tabular}{lrrrrr}
\toprule
Model & $K$ & RSS ($\times 10^9$) & AIC & AICc & $\Delta$AICc \\
\midrule
Logistic  & 217{,}781 & 3.66 & 652.2 & 653.0 & 11.5 \\
Gompertz  & 278{,}277 & 4.51 & 659.7 & 660.4 & 18.9 \\
Richards  & 183{,}618 & 2.44 & 640.2 & 641.5 &  0.0 \\
\bottomrule
\end{tabular}
\end{table}

The four-parameter Richards model provides the best AICc fit; the
logistic is second ($\Delta\mathrm{AICc} = 11.5$, moderate evidence
against logistic over Richards in conventional interpretation), and
Gompertz is clearly disfavoured ($\Delta\mathrm{AICc} = 18.9$).
The Richards shape parameter $\nu = 6.4$ indicates a noticeably
asymmetric approach to the carrying capacity relative to the symmetric
logistic, with a longer slow-growth tail before the inflection and a
sharper transition to the plateau. All three models agree qualitatively
on the main claim. Plant-science publication output is in the plateau
phase of sigmoidal growth, not in unconstrained exponential expansion.
We retain the logistic parameterisation in the main text for
interpretability ($K$ and inflection year $t_{\mathrm{mid}}$ have direct
meaning and are comparable to logistic fits reported elsewhere in the
metascience literature) and report the Richards comparison here so the
reader can assess the sensitivity of the plateau claim to model choice.

\begin{figure}[H]
  \centering
  \includegraphics[width=0.7\linewidth]{fig_growth_model_aic.pdf}
  \caption{\textbf{Growth-model comparison.} Observed annual plant-science
  publications (black points, 1990--2024) with logistic (blue),
  Gompertz (green dashed), and Richards (red dotted) fits. Legend
  reports $\Delta$AICc relative to the best-fitting model.}
  \label{fig:s-growth-aic}
\end{figure}

\suppnote{Climate-shock event study (pooled difference-in-differences)}{sec:climate-shocks}

To test whether major climate disasters produce measurable reactive
publication surges in the plant-science corpus, we applied a pooled
difference-in-differences event study to a pre-registered catalogue of
15 events (EM-DAT, NOAA Billion-Dollar Disasters, IPCC / UNFCCC policy
events; catalogue committed before analysis ran). For each event and
each affected-taxon category, we built monthly publication
time-series over $[t-36, +36]$ months, fit a pre-event linear trend
on $[-36, -3]$, and computed residuals for $[-3, +36]$. A matched
control series (publications on unrelated taxa over the same absolute
window, detrended identically) was subtracted. Event-specific effects
were pooled across events by averaging event-time residuals, with
95\% confidence intervals from 1{,}000-rep event-level bootstrap
resampling.

Taxon coverage was extended via an abstract-keyword path
(\texttt{monthly\_counts\_by\_keyword}) to resolve minor-crop events
(sorghum, millet, cassava, almond, grape, eucalyptus) and
topic-level policy events (IPCC AR5, Paris COP21, IPCC AR6) that
the organism classifier does not cover. With this extension, all 15
pre-registered events contributed non-empty treated series. The
pooled residual (Figure~\ref{fig:s-climate-shocks}) shows two
statistically significant response windows, a $+6$-month spike
(pooled mean above zero with CI $[$low $>0]$, likely reflecting
commentary and rapid-response reviews) and a larger $+18$-month
main spike (pooled mean $=+376.8$ papers/month, 95\% CI
$[+28.6, +773.8]$, excludes zero, consistent with full empirical
papers entering the literature 12--18 months after their source
experiments). Most intermediate lags show small or negative
residuals, indicating a burst rather than a smooth post-event
ramp. The 15-event coverage, the two-window response structure,
and the 2--3\% monthly-output redirection at the 18-month peak
together suggest that plant-science research output is measurably
responsive to climate pressure on a timescale consistent with
translational-research funding cycles.

\begin{figure}[H]
  \centering
  \includegraphics[width=0.7\linewidth]{fig_climate_shocks}
  \caption{\textbf{Pooled climate-event publication response.}
    Difference-in-differences residual (treated -- control) in
    publications per month, aligned to event time $t=0$ across
    $n=15$ events with non-empty treated series. Shaded band: 95\%
    bootstrap CI (1{,}000 event-level resamples).}
  \label{fig:s-climate-shocks}
\end{figure}

\suppnote{Paradigm-bridging authors --- pre-registered matched-impact analysis}{sec:bridge-scientists}

We tested whether authors who publish in both the
\textit{Arabidopsis} and crop paradigms (paradigm bridges) have
disproportionate citation impact relative to matched single-paradigm
peers. Authors active in 1990--2024 were classified as A-only
($\geq 3$ \textit{Arabidopsis}-tagged papers, zero crop), C-only
($\geq 3$ crop-tagged, zero \textit{Arabidopsis}), Bridge ($\geq 3$
in both), or Other (dropped). Each Bridge was matched to three
A-only and three C-only peers on first-paper year ($\pm 2$ years),
total paper count (within a factor of two), and primary
affiliation country (exact; relaxed to region when sparse).
Citations-per-paper were trimmed at the 99th percentile to guard
against superstar leverage.

Of 5{,}547 Bridge authors, 5{,}543 matched successfully (pool sizes:
A-only $= 19{,}299$; C-only $= 83{,}554$). Mean citations-per-paper
were: A-only $= 58.6$; Bridge $= 55.8$; C-only $= 30.1$
(Supplementary Fig.~\ref{fig:s-bridge-impact}).
Bridge-versus-C-only premium was $+85.5\%$ (95\% CI
$[+84.0, +87.0]$, permutation $p = 0.001$), clearing the
pre-registered $+50\%$ threshold. Bridge-versus-A-only was
$-4.7\%$ (95\% CI $[-5.6, -3.8]$), failing the pre-registered
threshold on the A-only comparison. The pre-registration therefore
rejects the original hypothesis ``Bridge authors are the citation
elite.'' The data support a different, complementary claim.
\textit{Arabidopsis}-specialist authors have a structural citation
advantage that bridging preserves (moving into crop work reduces
mean impact by less than 5\% relative to staying in
\textit{Arabidopsis}), while crop-only researchers have never had
that advantage (less than half the Bridge / A-only mean).
This pattern is consistent with the main-text D1 interpretation. If
\textit{Arabidopsis} literature functions as crop science's
foundational methodological substrate (Figure~\ref{fig:M9}),
\textit{Arabidopsis}-specialist authors attract citations precisely
because their work is the foundation that crop papers now cite.

\begin{figure}[H]
  \centering
  \includegraphics[width=0.55\linewidth]{fig_bridge_impact}
  \caption{\textbf{Mean citations-per-paper by paradigm class},
    99th-percentile trimmed, matched on first-paper year,
    productivity, and country. Permutation test (A-only vs Bridge):
    $p = 0.001$. Bars show group means; error bars are one standard
    error of the mean.}
  \label{fig:s-bridge-impact}
\end{figure}

\suppnote{Coauthor paradigm alignment over time (network-only)}{sec:coauthor-alignment}

To test whether the main-text $r_{\mathrm{inst}}$ peak at 2010--14
(Figure~\ref{fig:M9}a) corresponds to a structural inflection in the
field's social network, we built coauthor graphs for seven rolling
5-year windows (1990--1994 through 2020--2024) from the
\texttt{work\_authors} table and analysed them at the network level
rather than with community detection. Full Leiden community
detection on un-indexed DuckDB self-joins exceeded the sprint's
computational budget; a pre-aggregated coauthor-edge table is flagged
as follow-up work. We report two network-level metrics here,
degree-weighted paradigm alignment (the mean
\textit{Arabidopsis}-fraction of each author's portfolio weighted by
their coauthor-graph degree) and paradigm assortativity
(Newman's numeric assortativity coefficient on author portfolios).

Paradigm assortativity is high and nearly constant (0.82--0.88)
across all seven windows, indicating that plant-science coauthor
networks have been structurally paradigm-segregated throughout the
1990--2024 period --- this is not a new phenomenon. Paradigm
alignment, however, exhibits a clear trajectory, rising from
0.005 (1990--94) to a peak of 0.034 at 2010--14, then plateauing at
0.030 by 2020--24 (Figure~\ref{fig:s-paradigm-alignment}). The
2010--14 peak coincides with the independent $r_{\mathrm{inst}}$
peak in the main-text citation-flow analysis. Two methodologically
distinct measurements --- one on the citation graph, one on the
author--paradigm graph --- agreeing on the timing of the inflection
provides convergent evidence that the consolidation of
\textit{Arabidopsis} as crop science's methodological substrate
reached steady state in the mid-2010s.

\begin{figure}[H]
  \centering
  \includegraphics[width=0.7\linewidth]{fig_paradigm_alignment}
  \caption{\textbf{Paradigm alignment in the coauthor graph,
    1990--2024.} Degree-weighted mean \textit{Arabidopsis}-fraction of
    each author's portfolio, computed per 5-year window on a 10\%
    sampled coauthor edge set. Peaks at 2010--14, co-located with the
    main-text $r_{\mathrm{inst}}$ peak.}
  \label{fig:s-paradigm-alignment}
\end{figure}

\suppnote{Canon churn --- the plant-science reading list, 1995--2024}{sec:canon-churn}

Does a 2024 plant-science paper read from the same canon as a 2004
or 1995 paper? To answer this, we computed the top-100 most-cited
references for publications in each of six 5-year windows
(1995--99 through 2020--24), restricting cited works to those with
$\mathrm{pub\_year} \leq \mathrm{window\_start}$ so we measure what
each era drew on rather than concurrent citation swapping.

The canon is highly volatile. Only 30 of 100 top-100 references
from 1995--99 remain in the top-100 of 2020--24 --- over 70\% have
been displaced. Consecutive-window Jaccard overlap traces the
turnover rate, at 0.29 (1995--99 $\to$ 2000--04, lowest, likely
reflecting the genomics revolution), rising through 0.36, 0.40, to
a peak of 0.46 (2010--14 $\to$ 2015--19, the most stable period),
before dropping again to 0.37 at the 2015--19 $\to$ 2020--24
transition as deep-learning papers enter the canon
(Figure~\ref{fig:s-jaccard-over-time}). Kendall rank correlation
among shared references is 0.45--0.60 throughout, indicating not
just set turnover but within-set rank reshuffling. The reference
half-life distribution (Figure~\ref{fig:s-reference-half-life}) is
skewed, with 55 of 100 1995--99 top references appearing in that single
window only and six references appearing in all six windows (perennial
classics). Tables~\ref{tab:S11-risers} and \ref{tab:S11-fallen}
list the top-5 risers (new to 2020--24 top-100) and top-5 fallen
classics (in 1995--99 top-100, absent from 2020--24); three of the
top-five 2020--24 risers are deep-learning-for-agriculture papers,
and four of the top-five 1995--99 fallen classics are
mid-twentieth-century classical-methods papers.

\begin{table}[H]
\centering
\caption{Top-5 risers in the 2020--24 top-100 that did not appear in
any prior window's top-100.}
\label{tab:S11-risers}
\begin{tabular}{p{10cm}rr}
\toprule
Title & Year & Citations (2020--24) \\
\midrule
Using Deep Learning for Image-Based Plant Disease Detection & 2016 & 2{,}078 \\
Abiotic Stress Signaling and Responses in Plants & 2016 & 1{,}992 \\
Deep learning models for plant disease detection and diagnosis & 2018 & 1{,}547 \\
Shifting the limits in wheat research and breeding using a fully annotated reference genome & 2018 & 1{,}495 \\
Deep learning in agriculture: A survey & 2018 & 1{,}421 \\
\bottomrule
\end{tabular}
\end{table}

\begin{table}[H]
\centering
\caption{Top-5 fallen classics --- in 1995--99 top-100, absent from
2020--24 top-100.}
\label{tab:S11-fallen}
\begin{tabular}{p{10cm}rr}
\toprule
Title & Year & Citations (1995--99) \\
\midrule
A Two-Stage Technique for the In Vitro Digestion of Forage Crops & 1963 & 354 \\
A plant DNA minipreparation: Version II & 1983 & 348 \\
The estimation of protein degradability in the rumen from incubation measurements weighted according to rate of passage & 1979 & 345 \\
The Water-Culture Method for Growing Plants Without Soil & 1938 & 312 \\
A Manual for the Practical Study of Root-Nodule Bacteria & 1970 & 282 \\
\bottomrule
\end{tabular}
\end{table}

\begin{figure}[H]
  \centering
  \includegraphics[width=0.7\linewidth]{fig_jaccard_over_time}
  \caption{\textbf{Jaccard overlap of consecutive-window top-100
    reference sets, 1995--2024.} Lowest at the genomics-revolution
    transition (1995--99 $\to$ 2000--04); highest at the stable
    early-2010s (2010--14 $\to$ 2015--19); falls again at the
    deep-learning transition (2015--19 $\to$ 2020--24).}
  \label{fig:s-jaccard-over-time}
\end{figure}

\begin{figure}[H]
  \centering
  \includegraphics[width=0.7\linewidth]{fig_reference_half_life}
  \caption{\textbf{Reference half-life.} Number of 5-year windows
    a 1995--99 top-100 reference remains in any later top-100.
    Skewed distribution: 55 of 100 references appear only in 1995--99;
    six references persist across all six windows.}
  \label{fig:s-reference-half-life}
\end{figure}

\subsubsection*{Canon substitution --- is AI/ML displacing classical methods?}
\label{sec:canon-substitution}

The canon-churn analysis raises a natural follow-up. Is the rise of
deep-learning papers in the 2020--24 top-100 a \emph{substitution}
for the classical-methods papers that fell out of the canon, or are
the two trends independent? We tested this in two parts.

\paragraph*{Compositional trend (Part A).}
We classified each top-100 cited reference per window as
\emph{AI/ML} (abstract-keyword match on ``deep learning'',
``neural network'', ``machine learning'', ``convolutional'',
``computer vision'', etc.), \emph{Classical}
(``in vitro'', ``hydroponic'', ``gas chromatography'',
``mini preparation'', ``petri dish'', ``southern blot'', etc.),
or \emph{Other}. The AI/ML fraction is zero in every window from
1995--99 through 2015--19 and jumps to 4\% in 2020--24, while the
Classical fraction declines from 5\% (1995--99) to 1\% (2020--24).
Absolute fractions are low --- the canon remains dominated by
``Other'' (95\% throughout) --- but the first non-zero AI/ML entry
and the most-compressed Classical fraction coincide in the same
window (Figure~\ref{fig:s-composition-over-time}), consistent with
an emerging compositional pivot rather than drift.

\paragraph*{Pairwise substitution (Part B).}
For each of the top-5 2020--24 risers (three deep-learning-for-agriculture
papers + two stress-physiology reviews) and each of the top-5
1995--99 fallen classics (four mid-20th-century methods classics + one
microbiology manual), we computed TF-IDF cosine similarity of their
abstracts and the Pearson correlation between the fallen classic's
annual citation rate and the riser's annual citation rate over
overlapping years. Of the 25 (fallen~$\times$ riser) pairs, 7 show
$|r| > 0.5$ with opposite signs (fallen declining while riser rises);
the top three anti-correlated pairs are listed in
Table~\ref{tab:S11-substitution} and the strongest is visualised in
Figure~\ref{fig:s-top-substitution-pair}.

\begin{table}[H]
\centering
\caption{Top anti-correlated (fallen~$\times$ riser) pairs by Pearson~$r$
of citation-rate trajectories.}
\label{tab:S11-substitution}
\begin{tabular}{p{6cm}p{6cm}rr}
\toprule
Fallen classic & AI/ML riser & $r$ & $p$ \\
\midrule
Manual for practical study of root-nodule bacteria (1970) &
Using Deep Learning for Image-Based Plant Disease Detection (2016) &
$-0.853$ & $0.003$ \\
Manual for practical study of root-nodule bacteria (1970) &
Deep learning models for plant disease detection and diagnosis (2018) &
$-0.805$ & $0.016$ \\
The Water-Culture Method for Growing Plants Without Soil (1938) &
Using Deep Learning for Image-Based Plant Disease Detection (2016) &
$-0.712$ & $0.032$ \\
\bottomrule
\end{tabular}
\end{table}

We caution against reading these anti-correlations as demonstrated
causal displacement. TF-IDF abstract similarity across the 25 pairs is
close to zero (maximum $0.024$), and the 1970 microbiology manual and the
2016 CNN paper share essentially no vocabulary. The correlations
therefore capture temporal \emph{coincidence} --- these particular
classical references declined in exactly the years these particular
deep-learning papers rose --- without establishing that the
deep-learning papers displaced the classical ones in any specific
reader's reference list. A proper causal-substitution test would
require abstract-embedding-based topical proximity (e.g.\ SPECTER2
cosine distance) and within-topic citation-cohort analysis; we flag
this as follow-up work in Supplementary Note~\ref{sec:canon-churn}.

\begin{figure}[H]
  \centering
  \includegraphics[width=0.7\linewidth]{fig_composition_over_time}
  \caption{\textbf{Top-100 canon composition, 1995--2024.} Fraction
    of top-100 cited references classified as AI/ML (abstract keywords
    including ``deep learning'', ``neural network''), Classical
    (``in vitro'', ``hydroponic'', ``chromatography''), and Other,
    per 5-year window. AI/ML crosses Classical in the 2020--24
    window for the first time.}
  \label{fig:s-composition-over-time}
\end{figure}

\begin{figure}[H]
  \centering
  \includegraphics[width=0.7\linewidth]{fig_top_substitution_pair}
  \caption{\textbf{Top anti-correlated (fallen $\times$ riser) citation
    trajectories.} Annual citation rate of the root-nodule bacteria
    manual (1970; blue, left axis) and the deep-learning plant-disease
    detection paper (2016; red, right axis). Pearson correlation of
    the two series: $r = -0.853$ ($p = 0.003$). TF-IDF similarity of
    the two abstracts: $\approx 0$ (the papers share no vocabulary),
    so the anti-correlation reflects temporal coincidence rather than
    topical displacement.}
  \label{fig:s-top-substitution-pair}
\end{figure}
